\chapter{The Words of the Salam Language}
\label{app:keywords}

The words in the first table and in the section ``Other words of the
language'' are reserved for the Salam language itself and cannot be used as
the name of a variable or a function; neither on their own nor as part of a
multi-word name (the one exception is \code{main}, which is not a reserved
word, but since it is the name of the function where the program starts, a
variable of the same name cannot be made either). The first table lists the
words used in this book, and the section ``Other words of the language''
lists the words you will meet in the next books.

\section{Words used in this book}

\begin{center}
\begin{tabular}{lp{8.5cm}l}
\toprule
Word & Meaning and use & Chapter \\
\midrule
\code{func} & defines a function; \code{func main} is where the program starts & \ref{ch:first}, \ref{ch:functions} \\
\code{main} & the name of the function where the program starts (not reserved) & \ref{ch:first} \\
\code{end} & the end of every block (function, condition, loop) & \ref{ch:first} \\
\code{println} & writes to the output and moves to the next line & \ref{ch:first} \\
\code{print} & writes to the output without moving to the next line & \ref{ch:first} \\
\code{mut} & declares a variable whose value changes & \ref{ch:variables} \\
\code{const} & declares a constant & \ref{ch:variables} \\
\code{true} & the boolean value true & \ref{ch:types} \\
\code{false} & the boolean value false & \ref{ch:types} \\
\code{as} & type conversion & \ref{ch:types} \\
\bottomrule
\end{tabular}
\end{center}

\begin{center}
\begin{tabular}{lp{8.5cm}l}
\toprule
Word & Meaning and use & Chapter \\
\midrule
\code{and} & logical ``and'' & \ref{ch:operators} \\
\code{or} & logical ``or'' & \ref{ch:operators} \\
\code{not} & logical ``not'' (negation) & \ref{ch:operators} \\
\code{if} & a condition & \ref{ch:conditions} \\
\code{else} & the other path of a condition; with or without a condition & \ref{ch:conditions} \\
\code{match} & compares one value against several cases & \ref{ch:conditions} \\
\code{repeat} & a loop with a fixed count or a range & \ref{ch:loops} \\
\code{to} & the end of the range in \code{repeat} & \ref{ch:loops} \\
\code{by} & the size of the step in \code{repeat} (\code{repeat 1 to 10 by 2 in k}) & \ref{ch:loops} \\
\code{until} & a loop that runs while a condition holds & \ref{ch:loops} \\
\code{break} & ends a loop & \ref{ch:loops} \\
\code{continue} & skips to the next iteration of a loop & \ref{ch:loops} \\
\code{ret} & hands a result back from a function & \ref{ch:functions} \\
\code{each} & walks through the elements of an array & \ref{ch:arrays} \\
\code{in} & names the array in an \code{each} loop and the counter in \code{repeat} & \ref{ch:arrays}, \ref{ch:loops} \\
\code{eq} & equality comparison, the same as \code{==} & \ref{ch:operators} \\
\code{neq} & inequality comparison, the same as \code{!=} & \ref{ch:operators} \\
\code{lt} & less than, the same as \code{<} & \ref{ch:operators} \\
\code{gt} & greater than, the same as \code{>} & \ref{ch:operators} \\
\code{lte} & less than or equal, the same as \code{<=} & \ref{ch:operators} \\
\code{gte} & greater than or equal, the same as \code{>=} & \ref{ch:operators} \\
\bottomrule
\end{tabular}
\end{center}

\section{Type names}

\begin{center}
\begin{tabular}{ll}
\toprule
Name & Type \\
\midrule
\code{int} & integer \\
\code{f64} & floating-point number \\
\code{str} & text \\
\code{bool} & true or false \\
\bottomrule
\end{tabular}
\end{center}

Salam has other types as well (such as \code{i64} for very large integers,
\code{char} for a single one-byte character such as an English letter, and
\code{uchar} for a single Unicode character such as an accented or Persian
letter) which were not needed in this book. Type names are not reserved.

\section{Other words of the language}

These words are reserved too and must not be used in names:
\code{struct}, \code{enum}, \code{type}, \code{import}, \code{package},
\code{null}, \code{this}, \code{layout}, \code{on}, \code{pub},
\code{interface}, \code{impl}, \code{component}, \code{include},
\code{export}, \code{extern}, \code{defer}, \code{input}, \code{printerr},
\code{printerrln}, \code{pure}, \code{inline}, \code{noinline},
\code{noret}, \code{deprecated} and \code{switch}.

The words \code{link}, \code{static}, \code{dynamic} and \code{framework}
are not reserved; they have a special meaning only in the instruction that
links external libraries into a program, and elsewhere you may use them in
names.

\begin{note}
If a name you have chosen gets a ``reserved word'' error, or an odd syntax
error such as ``expected an expression'', one of its words may be reserved.
The simplest remedy is to change it a little: instead of \code{package}
write \code{parcel} or \code{box}; instead of \code{items in cart} write
\code{cart items}.
\end{note}

\section{Capabilities of strings and arrays}

These are not reserved words, but they attach to strings and arrays with a
dot, and they were used in this book:

\begin{center}
\begin{tabular}{ll}
\toprule
Capability & Job \\
\midrule
\code{len()} & the number of compartments of a string (chapter \ref{ch:strings}) or of elements of an array \\
\code{substr(start, count)} & cuts a piece out of a string \\
\code{split(separator)} & splits a string into a list of pieces \\
\code{get(index)} & reads one element of the list produced by \code{split} \\
\code{trim()} & removes the spaces at the start and end \\
\code{find(text)} & the position of a text inside a string, or \code{-1} \\
\code{to\_int()} & converts a string to an integer \\
\code{to\_float()} & converts a string to a floating-point number \\
\bottomrule
\end{tabular}
\end{center}
